\section{Diversity and Quality Control in Synthetic Data} \label{sec:25-background-synthetic}

\subsection{Generative Adversarial Networks}
\label{sec:25-background-synthetic-gans}

\begin{foldable}
  Generative adversarial networks~\cite{goodfellow2014generative} (GANs) frame image synthesis as a two-player game.
  A generator $G: \mathcal{Z} \rightarrow \mathcal{X}$ maps samples from a low-dimensional noise prior $\mathbf{z} \sim p_z$ to images in the data space $\mathcal{X}$.
  A discriminator $D: \mathcal{X} \rightarrow [0,1]$ attempts to distinguish real images from generated ones.
  The two networks are trained adversarially:
  \begin{equation}
    \min_{G} \max_{D} \; \mathbb{E}_{\mathbf{x} \sim p_{\text{data}}}[\log D(\mathbf{x})]
    + \mathbb{E}_{\mathbf{z} \sim p_z}[\log(1 - D(G(\mathbf{z})))].
    \label{eq:gan-minimax}
  \end{equation}
  At equilibrium, $G$ produces samples from $p_{\text{data}}$ and $D$ cannot distinguish real from generated images.
  In practice, this equilibrium is rarely achieved cleanly; two failure modes dominate.
\end{foldable}

\begin{foldable}
  \textbf{Vanishing gradients} occur when the discriminator becomes too strong early in training and assigns near-zero probability to all generated images.
  The generator gradient $\nabla_{\theta_G} \log(1 - D(G(\mathbf{z})))$ saturates, and learning stalls.
  The standard fix --- using $\nabla_{\theta_G} [-\log D(G(\mathbf{z}))]$ instead --- partially alleviates saturation but does not address the underlying distributional mismatch.
\end{foldable}

\begin{foldable}
  \textbf{Mode collapse} occurs when the generator maps a large region of noise space to a small region of the data space, producing diverse noise vectors that yield similar or identical outputs.
  The discriminator is fooled by a narrow set of ``easy'' images, and the generator has no incentive to explore the full data distribution.
  Mode collapse is particularly damaging in the KD context because a collapsed generator cannot produce the diversity of training queries needed to extract the full knowledge of the teacher.
\end{foldable}

\begin{foldable}
  The Wasserstein GAN~\cite{arjovsky2017wasserstein} (WGAN) replaces the Jensen-Shannon divergence implicit in the original GAN objective with the Wasserstein-1 (Earth Mover) distance, enforced by clipping discriminator weights to a compact set.
  The Wasserstein distance is defined even when the real and generated distributions have non-overlapping support, providing meaningful gradients throughout training and largely eliminating the vanishing gradient problem.
  WGAN-GP~\cite{gulrajani2017improved} replaces weight clipping with a gradient penalty, penalising the discriminator when its gradient norm deviates from one along interpolations between real and generated samples.
  This soft Lipschitz constraint produces more stable training and better mode coverage than clipping, and WGAN-GP has become a standard component in data-free distillation generators.
\end{foldable}

\begin{foldable}
  A key point for what follows: WGAN-GP improves training stability and reduces the severity of mode collapse, but it does not eliminate it in low-data regimes.
  When the generator is trained on a small support --- as in the few-shot and data-free settings of this thesis --- it can still learn a stable but narrow distribution that scores well on the
  Wasserstein objective while missing large regions of the target space.
  Stability is a necessary but not sufficient condition for diversity.
\end{foldable}

\subsection{The Diversity Problem in Knowledge Distillation}
\label{sec:25-background-synthetic-problem}

\begin{foldable}
  Mode collapse can be understood as a distributional coverage failure: the generated distribution $p_G$ approximates some subset of $p_{\text{data}}$ faithfully but assigns zero or negligible mass to a large portion of the real distribution.
  This is not simply a quality failure --- individual generated images may be indistinguishable from real ones --- it is a \emph{coverage} failure.
\end{foldable}

\begin{foldable}
  In the standard KD setting, mode collapse in the generator has a direct and measurable consequence.
  A student trained on a collapsed synthetic set learns to approximate the teacher's predictions on the modes present in that set.
  On held-out data that covers the full class distribution, the student's performance degrades in proportion to the fraction of the distribution that was absent from its training set.
  The binding constraint on student accuracy is not the fidelity of individual generated images but the distributional coverage of the generated set as a whole.
\end{foldable}

\begin{foldable}
  This framing is broader than the image KD context.
  In dataset distillation (Ch.~\ref{ch:50-research03}), the analogous failure is a corpus whose selected examples cluster around the high-frequency training instances and fail to cover low-frequency but semantically distinct examples.
  In both cases, the question is the same: does the compact artefact --- synthetic image set or selected text corpus --- cover the distribution it is meant to represent?
  Distributional coverage is therefore the unifying concern across both tracks of the thesis.
\end{foldable}

\subsection{Evaluation Metrics for Diversity and Generation Quality}
\label{sec:25-background-synthetic-metrics}

\begin{foldable}
  Standard image generation metrics conflate fidelity and diversity in ways that make them insensitive to coverage failures.
\end{foldable}

\begin{foldable}
  \textbf{\ac{IS}}~\cite{salimans2016improved} measures two properties jointly: the conditional label distribution $p(y \mid \mathbf{x})$ of generated images should be peaked (high fidelity, low entropy per image), and the marginal distribution $p(y) = \int p(y \mid \mathbf{x}) p_G(\mathbf{x}) \, \mathrm{d}\mathbf{x}$ should be spread (high diversity across images).
  \ac{IS} is computed as $\exp(\mathbb{E}_{\mathbf{x} \sim p_G}[\mathrm{KL}(p(y|\mathbf{x}) \| p(y))])$.
  Despite its widespread use, \ac{IS} does not measure \emph{intra-class} diversity and does not detect mode collapse within individual classes.
\end{foldable}

\begin{foldable}
  \textbf{\ac{FID}}~\cite{heusel2017gans} models both real and generated distributions as Gaussians in Inception feature space and measures the Fréchet distance between them:
  \begin{equation}
    \mathrm{FID} = \|\mu_r - \mu_g\|^2 + \mathrm{Tr}\!\left(\Sigma_r + \Sigma_g
    - 2(\Sigma_r \Sigma_g)^{1/2}\right),
    \label{eq:fid}
  \end{equation}
  where $(\mu_r, \Sigma_r)$ and $(\mu_g, \Sigma_g)$ are the empirical means and covariances of real and generated features.
  \ac{FID} is sensitive to both the mean (fidelity) and covariance (diversity) of the generated distribution, making it a better overall metric than \ac{IS}.
  However, it still aggregates fidelity and diversity into a single number and does not separately report how much of the real distribution is covered.
\end{foldable}

\begin{foldable}
  \textbf{Density and Coverage}~\cite{naeem2020reliable} decompose the distributional similarity into precision-like and recall-like components.
  Density measures how reliably generated samples fall near real samples (fidelity); Coverage measures what fraction of the real distribution is represented by the generated samples (recall).
  Formally, Coverage is computed as the proportion of real samples whose neighbourhood in feature space contains at least one generated sample.
  Coverage is the primary metric for the diversity problem: a set of generated images can have high density (individually high-quality) and low coverage (distributed over only a few modes).
  All three research chapters report Coverage (or equivalent recall metrics) in addition to FID, enabling a direct assessment of distributional coverage.
\end{foldable}

\begin{foldable}
  \textbf{Downstream classification accuracy} is the ultimate evaluation criterion throughout the thesis: the purpose of both model KD and dataset DD is to train a student that performs well on the task, not merely to produce images that score well on generative metrics.
  Generative metrics are used as diagnostic tools; accuracy is the primary measure of success.
\end{foldable}

\begin{foldable}
  The three metrics above form a diagnostic hierarchy rather than a flat set of alternatives.
  Coverage establishes whether the generated set spans the target distribution --- it is the necessary condition for avoiding silent failure modes in which individual samples are high-quality but entire regions of the distribution are absent.
  \ac{FID} then quantifies the aggregate distributional gap: even when coverage is adequate, a large \ac{FID} signals that the generator has distorted the shape of the distribution in some way.
  Downstream accuracy is the bottom line that both metrics serve --- it is the empirical verdict on whether the coverage and fidelity achieved at the generative level are sufficient for the student learning task.
  Crucially, these metrics also reveal a limitation of purely generative objectives.
  A generator trained to minimise \ac{FID} optimises for aggregate distributional similarity; it may achieve a low \ac{FID} by accurately capturing the high-density modes of the target distribution while leaving low-density regions unrepresented, a failure that is invisible in the single aggregate score but that Coverage would detect.
  Improving \ac{FID} therefore does not guarantee improving Coverage, and vice versa.
  This dissociation points to a need that evaluation alone cannot satisfy: a mechanism that actively enforces distributional breadth during generation, before the distribution of the generated set is ever measured.
  The contrastive learning framework described in the following subsection provides exactly this mechanism.
\end{foldable}

\subsection{Contrastive Learning as Diversity Enforcement}
\label{sec:25-background-synthetic-contrastive}

\begin{foldable}
  Contrastive self-supervised learning provides a principled mechanism for enforcing representational diversity without requiring labels.
  The core objective is to learn an encoder $f_\theta: \mathcal{X} \rightarrow \mathbb{R}^d$ such that two augmented views of the same image (a positive pair) map to nearby representations, while views of different images (negative pairs) map to distant representations.
  This \emph{attraction-repulsion} structure directly prevents representation collapse: if all images were mapped to the same embedding, the repulsion term would incur a large loss.
\end{foldable}

\begin{foldable}
  SimCLR~\cite{chen2020simple} implements this objective using the Normalised Temperature-scaled Cross-Entropy (NT-Xent) loss over a large batch of positive and negative pairs.
  For a positive pair $(i, j)$ of views drawn from the same image, the loss is:
  \begin{equation}
    \ell_{i,j} = -\log \frac{\exp(\mathrm{sim}(\mathbf{z}_i, \mathbf{z}_j) / \tau)}
    {\sum_{k \neq i} \exp(\mathrm{sim}(\mathbf{z}_i, \mathbf{z}_k) / \tau)},
    \label{eq:ntxent}
  \end{equation}
  where $\mathbf{z}_i = g(f_\theta(\tilde{\mathbf{x}}_i))$ is the projection of augmented view $\tilde{\mathbf{x}}_i$, $\mathrm{sim}(\cdot, \cdot)$ is cosine similarity, and $\tau$ is a temperature parameter.
  The denominator sums over all other images in the batch, acting as the negative set.
\end{foldable}

\begin{foldable}
  MoCo~\cite{he2020momentum} addresses the large-batch requirement of SimCLR by maintaining a memory queue of negative embeddings encoded by a momentum-updated copy of the encoder.
  This decouples the negative set size from the batch size, enabling effective contrastive training on hardware with limited memory.
\end{foldable}

\begin{foldable}
  Barlow Twins~\cite{zbontar2021barlow} takes a different approach: instead of explicitly comparing instances, it encourages the cross-correlation matrix between the embeddings of two augmented views to be close to the identity matrix.
  This redundancy-reduction objective drives each embedding dimension to be independently informative, avoiding the need to manage negative pairs explicitly.
\end{foldable}

\begin{foldable}
  The relevance of contrastive objectives to the thesis is geometric.
  The repulsion term in any contrastive loss drives embeddings apart in the representation space, explicitly maximising the separation between different instances.
  A generator whose outputs are diverse in embedding space is, by construction, covering a larger region of the target distribution.
  Chapter~\ref{ch:40-research02} exploits this property directly: the Contrast phase of DIPKD applies a contrastive objective to the primer student's features over the synthetic image priors, expanding semantic coverage before the main distillation phase begins.
\end{foldable}
